\begin{notechapter}
\chapter{Dropout}\label{ch:dropout}
\chaptersubtitle{Random masks at training time, and what to do with them at test time}

\begin{chapterintro}
Dropout sets a random subset of the units of a layer to zero at each training step. It is
one of the most effective regularizers for large networks, and also one whose test-time
behaviour is easy to get wrong. This chapter defines it precisely (with one convention for
the probability), explains why the network used at test time is an approximation, and
when that approximation is exact. You will learn
\begin{itemize}[nosep]
  \item standard and inverted dropout, and their means and variances;
  \item the ensemble interpretation, the weight-scaling rule, when it is exact (linear
        layers, and a softmax output in the sense of a geometric mean) and when it is not;
  \item that dropout on the inputs of a linear regression is, in expectation, a ridge penalty;
  \item Monte Carlo dropout and the variational interpretation;
  \item practical values, variants for convolutional layers, and how to implement all of this
        in PyTorch without the usual mistakes.
\end{itemize}
\prerequisites{\cref{ch:feedforward-networks}; ridge regression and bias--variance
(\cref{ch:erm}); Bayesian inference and the ELBO for the last section
(\cref{ch:bayesian-reasoning,ch:em-algorithm}).}
\end{chapterintro}

\section{Definition}\label{dropout:sec:definition}

\paragraph{Convention.} In this chapter $p$ is the \emph{drop} probability and $q = 1 - p$ the \emph{keep}
probability. This is PyTorch's convention (\texttt{nn.Dropout(p)}); the original papers use $p$
for the keep probability.

\begin{definition}[Dropout]\label{dropout:def:dropout}
Let $a\in\R^n$ be the activations of a layer. At each training step draw a mask $m\in\{0,1\}^n$ with
independent coordinates $m_j\sim\mathrm{Bernoulli}(q)$, independently for each example, and output
\[
  \text{standard:}\quad\tilde a = m\odot a,\qquad\qquad\text{inverted:}\quad\tilde a = \frac1q\,m\odot a .
\]
At test time standard dropout outputs $qa$ and inverted dropout outputs $a$.
\end{definition}

\begin{center}
\begin{tabular}{lcc}
\toprule
& training & test \\
\midrule
standard dropout & $m\odot a$ & $q\,a$ \\
inverted dropout & $m\odot a/q$ & $a$ \\
\bottomrule
\end{tabular}
\end{center}

\begin{proposition}[Moments]\label{dropout:prop:moments}
For fixed $a$, inverted dropout gives $\E[\tilde a_j] = a_j$ and $\Var(\tilde a_j) = \frac pqa_j^2$; standard dropout gives
$\E[\tilde a_j] = qa_j$ and $\Var(\tilde a_j) = pq\,a_j^2$.
\end{proposition}

\begin{proof}
$\E[m_j] = q$ and $\Var(m_j) = q(1 - q) = pq$; for inverted dropout divide the mean by $q$ and the variance by $q^2$.
\end{proof}

Both versions make the test-time output equal to the training-time \emph{mean} of each unit.
They define the same model up to a rescaling of the next layer's weights by $q$; inverted
dropout, which leaves the test-time network untouched, is what libraries implement. For
$p = 0.5$ the training-time output of a kept unit is $2a_j$, and its standard deviation equals its
mean: dropout is a strong multiplicative noise.

\section{Why it regularizes}\label{dropout:sec:why}

Srivastava, Hinton, Krizhevsky, Sutskever and Salakhutdinov (\emph{Dropout: A Simple Way to
Prevent Neural Networks from Overfitting}, JMLR, 2014), following Hinton et al.\ (2012),
give two intuitions.

\paragraph{Co-adaptation.} A unit cannot rely on the presence of particular other units,
which may be dropped. It is pushed to compute features that are useful in many contexts,
rather than to correct the specific mistakes of other units on the training set.

\paragraph{Ensembles.} Each mask selects a sub-network. With $n$ droppable units there are
$2^n$ sub-networks, all sharing the same weights, and each training step trains one of them on
one mini-batch. Dropout trains an exponentially large ensemble with parameter sharing; at
test time one would like to average their predictions (\cref{ch:ensemble-methods}).

\paragraph{An exact case: linear regression.} The regularizing effect can be computed when
the model is linear.

\begin{proposition}[Dropout on the inputs of a linear model]\label{dropout:prop:ridge}
Let $\hat y = \sum_jw_j\tilde x_j$ with inverted dropout on the inputs, $\tilde x_j = m_jx_j/q$. Then
\[
  \E_m\bigl[(y - \hat y)^2\bigr] = (y - w^\top x)^2 + \frac pq\sum_jw_j^2x_j^2 .
\]
Over a data set with design matrix $X$, the expected dropout loss is
$\norm{y - Xw}^2 + \frac pq\,w^\top\diag(X^\top X)\,w$: ridge regression with a penalty on each $w_j$ proportional to
$\sum_ix_{ij}^2$.
\end{proposition}

\begin{proof}
$\E[\hat y] = w^\top x$, and $\E[(y - \hat y)^2] = (y - \E\hat y)^2 + \Var(\hat y)$. The terms $w_j\tilde x_j$ are independent with variances
$w_j^2x_j^2\,p/q$ by \cref{dropout:prop:moments}. Summing over examples gives the matrix form.
\end{proof}

This is Section 9.1 of Srivastava et al.\ (2014); Wager, Wang and Liang (\emph{Dropout Training
as Adaptive Regularization}, NeurIPS, 2013) extend it, to second order, to generalized
linear models, where the penalty also depends on the predicted variances. The penalty grows
with $p/q$: from $0$ for $p = 0$ to $1$ for $p = \frac12$ and $\infty$ as $p\to1$. Unlike plain ridge, it is scaled per feature
by the feature's magnitude, which makes it invariant to rescaling the inputs.

\section{Inference}\label{dropout:sec:inference}

The ideal prediction averages over masks, $\E_m[f(x; m)]$, which has $2^n$ terms per dropout layer.
The \emph{weight-scaling rule} replaces each mask by its expectation, that is, uses the
deterministic network (inverted dropout) or the network with activations multiplied by $q$
(standard dropout). With several dropout layers, each mask is replaced independently.

\paragraph{Exact for linear maps.} If the network is linear in the dropped activations,
expectation and network commute and the rule is exact. This covers the last layer before
the output (the logits are linear in the last hidden activations).

\paragraph{Exact as a geometric mean for a softmax output.}

\begin{proposition}\label{dropout:prop:geometric}
Let $P_m(y\mid x) = \softmax\bigl(W(m\odot x)\bigr)_y$ with standard dropout on $x$. The weighted geometric mean of
these distributions over masks, renormalized,
\[
  \tilde P(y\mid x)\ \propto\ \prod_mP_m(y\mid x)^{\Prob(m)},
\]
equals $\softmax\bigl(W(qx)\bigr)_y$, the prediction of the weight-scaled network.
\end{proposition}

\begin{proof}
$\log P_m(y\mid x) = w_y^\top(m\odot x) - \log Z_m$, where $Z_m$ does not depend on $y$. Taking the expectation over $m$,
$\sum_m\Prob(m)\log P_m(y\mid x) = w_y^\top(qx) - c$ with $c$ independent of $y$. Exponentiating and normalizing
over $y$ gives $\softmax(W(qx))$.
\end{proof}

Hinton et al.\ (\emph{Improving Neural Networks by Preventing Co-adaptation of Feature
Detectors}, 2012) note this for one hidden layer with $p = \frac12$. It is the \emph{geometric}, not the
arithmetic, mean of the ensemble; and with hidden nonlinearities between dropout and
output, it no longer holds.

\begin{example}[The rule can fail]\label{dropout:ex:counter}
A ReLU unit receives two inputs $x = (1,1)$ with weights $w = (1,-1)$, with inverted dropout, $p = \frac12$:
$h = \mathrm{ReLU}(2m_1 - 2m_2)$. The four masks, each of probability $\frac14$, give $h = 0, 2, 0, 0$ for
$(m_1,m_2) = (1,1), (1,0), (0,1), (0,0)$, so $\E[h] = 0.5$. The deterministic network gives $\mathrm{ReLU}(1 - 1) = 0$. In
general $\E[\phi(Z)]\neq\phi(\E Z)$, and the rule only approximates the average over masks; in practice the
approximation is good enough that it is almost always used.
\end{example}

\paragraph{Monte Carlo dropout.} Keeping dropout active at test time and averaging $K$
stochastic forward passes,
\[
  \hat f(x) = \frac1K\sum_{k=1}^Kf(x; m^{(k)}),
\]
estimates $\E_m[f(x; m)]$ without bias, with a standard error decreasing like $1/\sqrt K$. The spread of the
$K$ outputs measures how much the prediction depends on the mask. Gal and Ghahramani
(\emph{Dropout as a Bayesian Approximation: Representing Model Uncertainty in Deep
Learning}, ICML, 2016) showed that training with dropout and weight decay can be read as
approximate variational inference (maximization of an ELBO, \cref{ch:em-algorithm}) in a
Bayesian network,
with an approximate posterior over weights given by the random masks. Under this reading, Monte
Carlo dropout approximates the posterior predictive distribution, and its spread is an
estimate of model uncertainty. The variational family is very restricted, so this
uncertainty is a rough approximation of the Bayesian one, not a sample from the true
posterior.

\section{In practice}\label{dropout:sec:practice}

\paragraph{Values.} $p = 0.5$ for large fully connected hidden layers and $p\approx0.2$ for inputs were the
original recommendations (Srivastava et al.\ 2014); transformers use $p = 0.1$ (Vaswani et al.,
\emph{Attention Is All You Need}, NeurIPS, 2017). Dropout is usually not applied to the
output layer. Since it slows convergence (noisier gradients), networks with dropout need
more epochs.

\paragraph{Variants.} In convolutional layers, neighbouring activations are strongly
correlated, so dropping single entries removes little information. \emph{Spatial dropout}
drops whole channels (Tompson, Goroshin, Jain, LeCun and Bregler, CVPR, 2015);
\emph{Cutout} masks a square patch of the input image (DeVries and Taylor, 2017);
\emph{DropConnect} drops individual weights instead of activations (Wan, Zeiler, Zhang, LeCun
and Fergus, ICML, 2013). Combined with batch normalization, dropout should come after the
normalization layers (\cref{ch:normalization}).

\paragraph{In PyTorch.} \texttt{nn.Dropout(p)} implements inverted dropout: in training mode
it multiplies by $m/(1 - p)$, in evaluation mode it is the identity.
\begin{lstlisting}[language=Python]
import torch
from torch import nn

model = nn.Sequential(nn.Linear(10, 20), nn.ReLU(),
                      nn.Dropout(p=0.5),        # drop probability 0.5
                      nn.Linear(20, 1))
x = torch.randn(16, 10)

model.train(); y1 = model(x)                    # random: masks are drawn
model.eval();  y2 = model(x)                    # deterministic: identity

def mc_dropout(model, x, K=50):                 # Monte Carlo dropout
    model.eval()                                # batch norm etc. in eval mode
    for mod in model.modules():
        if isinstance(mod, nn.Dropout):
            mod.train()                         # only dropout stays random
    with torch.no_grad():
        ys = torch.stack([model(x) for _ in range(K)])
    model.eval()
    return ys.mean(0), ys.std(0)
\end{lstlisting}

\begin{pitfall}[Monte Carlo dropout with \texttt{model.train()}]
Calling \texttt{model.train()} to keep dropout active at test time also puts batch
normalization layers in training mode: they then normalize with the statistics of the test
batch and update their running averages with test data. Activate the dropout modules only,
as in \texttt{mc\_dropout} above.
\end{pitfall}

\begin{chaptersummary}
\begin{itemize}[nosep]
  \item Dropout multiplies activations by independent $\mathrm{Bernoulli}(q)$ masks, $q = 1 - p$. Inverted dropout divides by
        $q$ during training, so that the test-time network is unchanged; the noise has variance $\frac pqa^2$.
  \item It trains $2^n$ weight-sharing sub-networks and discourages co-adaptation. For linear
        regression it is, in expectation, a ridge penalty $\frac pq\sum_jw_j^2\sum_ix_{ij}^2$.
  \item The weight-scaling rule is exact for linear maps, and for a softmax output in the sense
        of the normalized geometric mean; with hidden nonlinearities it is an approximation.
  \item Monte Carlo dropout averages stochastic passes; it estimates the mean over masks and,
        under the variational interpretation, a crude model uncertainty. Enable only the dropout
        modules when doing it.
\end{itemize}
\end{chaptersummary}

\section*{Exercises}
\addcontentsline{toc}{section}{Exercises}

\begin{exercise}\label{dropout:exr:1}
With inverted dropout and $p = 0.2$, compute the mean and standard deviation of $\tilde a_j$ for $a_j = 1$. For
which $p$ does the standard deviation equal the mean?
\solutionlink{dropout:exr:1}
\end{exercise}

\begin{exercise}\label{dropout:exr:2}
Two features $x_1, x_2$ with $\sum_ix_{i1}^2 = 100$ and $\sum_ix_{i2}^2 = 1$. With $p = 0.5$ on the inputs, write the expected
dropout objective of \cref{dropout:prop:ridge}. If the second feature is multiplied by $10$ and $w_2$ divided
by $10$, does the penalty change? Compare with ordinary ridge.
\solutionlink{dropout:exr:2}
\end{exercise}

\begin{exercise}\label{dropout:exr:3}
In \cref{dropout:prop:geometric}, take two classes, $x = (2,1)$, $W = \begin{bmatrix}1&0\\0&1\end{bmatrix}$, $q = \frac12$. Compute the four
distributions $P_m$, their arithmetic mean, their normalized geometric mean, and
$\softmax(W(qx))$.
\solutionlink{dropout:exr:3}
\end{exercise}

\begin{exercise}\label{dropout:exr:4}
In \cref{dropout:ex:counter}, compute the variance of $h$ over the masks, and the number $K$ of Monte
Carlo passes needed for a standard error below $0.05$.
\solutionlink{dropout:exr:4}
\end{exercise}

\begin{exercise}\label{dropout:exr:5}
Show that standard dropout followed at test time by multiplication by $q$ and inverted
dropout define the same family of predictors, by relating the weights of the next layer
in the two versions.
\solutionlink{dropout:exr:5}
\end{exercise}
\end{notechapter}
